\begin{mistake}{2026-04-09}{37}

\begin{problemcontent}
设 \( y_1, y_2 \) 是 \( y'' + py' + qy = f(x) \) 的解，\( y_3, y_4 \) 是 \( y'' + py' + qy = g(x) \) 的解。以下哪个是 \( y'' + py' + qy = f(x) - g(x) \) 的解？\\
(A) \( y_1 - 2y_2 + y_3 - 2y_4 \)\\
(B) \( 2y_1 - y_2 + y_3 - 2y_4 \)\\
(C) \( 2y_1 - y_2 + 2y_3 - y_4 \)\\
(D) \( y_1 - 2y_2 + 2y_3 - y_4 \)
\end{problemcontent}

\begin{solutioncontent}
设微分算子 \( L[y] = y'' + py' + qy \)。由题意知 \( L[y_1]=L[y_2]=f(x) \)，\( L[y_3]=L[y_4]=g(x) \)。

令特解形式为 \( y = C_1y_1 + C_2y_2 + C_3y_3 + C_4y_4 \)。将其代入方程左侧，利用线性叠加原理有：
\[
\begin{aligned}
L[y] &= C_1L[y_1] + C_2L[y_2] + C_3L[y_3] + C_4L[y_4] \\
&= (C_1+C_2)f(x) + (C_3+C_4)g(x)
\end{aligned}
\]
若要使其为 \( f(x) - g(x) \) 的解，必须满足系数约束条件：
\( C_1 + C_2 = 1 \) 且 \( C_3 + C_4 = -1 \)。

逐一检验选项：\\
(A) \( C_1+C_2 = -1 \)，不符。\\
(B) \( C_1+C_2 = 2-1=1 \)，\( C_3+C_4 = 1-2=-1 \)，符合。\\
(C) \( C_3+C_4 = 2-1=1 \)，不符。\\
(D) \( C_1+C_2 = 1-2=-1 \)，不符。\\
故选 (B)。
\end{solutioncontent}

\mistakeknowledge{
\begin{itemize}
 \item \textbf{核心定理}：非齐次线性微分方程的叠加原理 \( L[C_1y_1 + C_2y_2] = C_1L[y_1] + C_2L[y_2] \)。
 \item \textbf{易错点}：目标方程右侧 \( g(x) \) 前为负号，极易遗漏导致错选 \( C_3+C_4=1 \) 的选项。
 \item \textbf{关联知识}：若 \( y_1, y_2 \) 同为某非齐次方程的解，则 \( y_1 - y_2 \) 为对应齐次方程的解。
\end{itemize}}

\end{mistake}
